%% Supplementary Material — KNOSYS-D-26-21490
%% Knowledge-Based Systems, Elsevier

\documentclass[a4paper,11pt]{article}
\usepackage[T1]{fontenc}
\usepackage[utf8]{inputenc}
\usepackage{geometry}
\geometry{margin=2.5cm}
\usepackage{booktabs}
\usepackage{tabularx}
\usepackage{url}
\usepackage{microtype}
\usepackage{authblk}
\usepackage{hyperref}
\hypersetup{colorlinks=true, linkcolor=blue, urlcolor=blue}

\title{\textbf{Supplementary Material}\\[4pt]
\large Effective Training Pressure Gates Recursive Knowledge\\
Degradation in LLMs: A Multi-Axis Dose-Response Study}

\author[1,2]{Julio Leite Azancort Neto%
  \thanks{Corresponding author. E-mail: \href{mailto:jneto@ufpa.br}{jneto@ufpa.br}.
  ORCID: 0000-0003-2866-5445}}
\author[1,2]{Carlos Andr\'{e} de Mattos Teixeira%
  \thanks{ORCID: 0000-0003-1902-4347}}
\author[3]{Andr\'{e} Carlos Ponce de Leon Ferreira de Carvalho%
  \thanks{ORCID: 0000-0002-1907-5913}}
\author[1,2]{Carlos Renato Lisboa Franc\^{e}s%
  \thanks{ORCID: 0000-0003-0305-7662}}

\affil[1]{Amazon Sustainability Institute with Science and Innovation
  (iSACI), Bel\'{e}m, Par\'{a}, Brazil}
\affil[2]{Programa de P\'{o}s-Gradua\c{c}\~{a}o em Engenharia
  El\'{e}trica e de Computa\c{c}\~{a}o (PPGEE), Universidade Federal
  do Par\'{a} (UFPA), Bel\'{e}m, Par\'{a}, Brazil}
\affil[3]{Instituto de Ci\^{e}ncias Matem\'{a}ticas e de
  Computa\c{c}\~{a}o (ICMC), Universidade de S\~{a}o Paulo (USP),
  S\~{a}o Carlos, SP, Brazil}

\date{}

\setcounter{table}{0}
\renewcommand{\thetable}{S\arabic{table}}

\begin{document}

\maketitle

\begin{abstract}
\noindent This supplementary document provides the base-model checkpoint
identifiers (Table~S1) used in the experiments reported in the main
manuscript.
\end{abstract}

% ─────────────────────────────────────────────────────────────────────────────
\section*{Base-Model Checkpoint Identifiers}
% ─────────────────────────────────────────────────────────────────────────────

All three backbones were loaded from the Hugging Face Hub using
\texttt{AutoModelForCausalLM.from\_pretrained(repo\_id)} without an
explicit \texttt{revision=} argument. The SHA-1 hashes in Table~S1
identify the snapshots present in the local cache at the time of each
experiment; weights can be retrieved via
\texttt{from\_pretrained(repo\_id, revision=sha1)}.

\textit{Provenance note for Gemma~4~E2B~IT.}\;
Snapshot \texttt{70af34e} was downloaded on 2026-06-26; the Gemma~4
experiments ran on 2026-06-27, when this was the only snapshot in the
local cache. A second snapshot (\texttt{3e22461}) entered the cache on
2026-09-30 and was not used in any reported experiment.

\begin{table}[ht]
\caption{HuggingFace Hub repository identifiers and cache-snapshot
SHA-1 hashes for the three base-model checkpoints used in this study.
$^{\dagger}$See provenance note above.}
\label{tab:s1}
\footnotesize
\begin{tabularx}{\textwidth}{@{}l l X@{}}
\toprule
\textbf{Backbone} & \textbf{SHA-1 (40 chars)} & \textbf{Repository} \\
\midrule
Qwen~2.5-1.5B-Instruct
  & \texttt{989aa7980e4cf806f80c7fef2b1adb7bc71aa306}
  & \texttt{Qwen/Qwen2.5-1.5B-Instruct} \\[4pt]
Gemma~3~1B~IT
  & \texttt{dcc83ea841ab6100d6b47a070329e1ba4cf78752}
  & \texttt{google/gemma-3-1b-it} \\[4pt]
Gemma~4~E2B~IT\,$^{\dagger}$
  & \texttt{70af34e20bd4b7a91f0de6b22675850c43922a03}
  & \texttt{google/gemma-4-e2b-it} \\
\bottomrule
\end{tabularx}
\end{table}

\noindent These identifiers are also listed in
\texttt{configs/checkpoints.yaml} of the replication package
(Zenodo DOI: \url{10.5281/zenodo.23146342}).

\end{document}
